\section{Discussion and open questions}\label{sec:discussion}

The empirical results of Section~\ref{sec:results} show that the
SMC${}^2$ filter and controller, composed in receding-horizon form,
produce a closed-loop policy that improves mean autonomic amplitude by
$+17\%$ over a sensible constant baseline on the FSA-v2 model, with no
fatigue-ceiling violations.

Earlier drafts of these notes flagged three analyses as outstanding:
Fisher information rank identifiability, Lyapunov stability, and a
linear--quadratic--Gaussian (LQG) approximation. The first two have
since been carried out and now form Sections~\ref{sec:fim}
and~\ref{sec:lyap}. The third remains as the principal open piece of
analytical work and is described in
Section~\ref{sec:discussion-lqg} below. Three smaller open questions
are then collected in Section~\ref{sec:discussion-other}.

\subsection{Resolved: Fisher information rank
identifiability}\label{sec:discussion-fim}

The pen-and-paper FIM analysis of Section~\ref{sec:fim} identified
three rank-deficient pairs in the original v2 parametrization,
$(\kappa_B,\varepsilon_A)$, $(\mu_F,\mu_{FF})$, and
$(\tau_F,\lambda_A)$. Each pair was shown to enter the one-day-window
likelihood only through a single linear combination, which is the
combinatorial signature of a rank-one rather than rank-two contribution
to the FIM. The Stage~G1 reparametrization rotates each pair into a
strongly-identified combination and a residual that vanishes at the
operating point of Definition~\ref{def:operating-point}; empirically
this restores full FIM rank and contracts the posterior standard
deviations on the residuals by roughly a factor of six (see
Table~\ref{tab:g1-evidence}).

A useful follow-up, identified in Remark~\ref{rem:fim-where}, is to
add a numerical FIM module under
\texttt{smc2fc/diagnostics/fim/} that recomputes the rank and
condition number directly from the JAX dynamics, so that any future
modification of the model triggers an automatic re-check.

\subsection{Resolved: Lyapunov stability of the
SDE}\label{sec:discussion-lyap}

Section~\ref{sec:lyap} carried out the stability analysis in two
parts. The deterministic skeleton, under a constant control
$\Phi \equiv \Phi_c$, has two branches: a sedentary $A=0$ branch and
an oscillatory $A^*_{\mathrm{osc}}(\Phi_c)$ branch. Local stability of
each is determined by the sign of the Stuart--Landau growth rate
$\mu(B^*_0,F^*_0)$. The Goldilocks profile of
Section~\ref{sec:lyap-goldilocks} -- $A^*_{\mathrm{osc}}$ rising
from a non-zero rest value at $\Phi=0$, peaking at the constant-load
optimum $\Phi_{\mathrm{opt}} \approx 1.99$, and collapsing at the
bifurcation locus $\Phi_c^{\mathrm{crit}} \approx 3.71$ as
overtraining suppresses $\mu$ -- explains both why the canonical
$\Phi=1$ baseline is suboptimal and why the SMC${}^2$ controller's
mean schedule clusters near $\Phi \approx 2$. For the full
SDE we exhibited a quadratic Lyapunov function
$V = \tfrac12(aB^2+bF^2+cA^2)$ that satisfies a Foster--Lyapunov drift
inequality (Theorem~\ref{thm:lyap}), and concluded exponential
ergodicity at rate
$\rho < \alpha = \min(2/\tau_B,1/\tau_F^{\mathrm{old}})$
(Theorem~\ref{thm:ergodic}). This delivers strong stability of the
controlled process on the operating domain
$\Xcal = [0,1]\times\R_+^2$ for any bounded control schedule.

Remark~\ref{rem:lyap-limits} noted that the constants in
Theorem~\ref{thm:lyap} are loose: they use the simplest quadratic
Lyapunov function and global bounds on the residual factors. A
sharper rate could be obtained from a piecewise Lyapunov function
that switches between the $\mu < 0$ and $\mu > 0$ regimes identified
in Proposition~\ref{prop:stab-zero}, or from a non-quadratic Lyapunov
function tuned to the geometry of the Stuart--Landau cubic.

\subsection{Outstanding: LQG / Riccati
exact-approximate alternative}\label{sec:discussion-lqg}

An earlier draft of these notes flagged the
linear--quadratic--Gaussian (LQG) / Riccati controller of
Section~\ref{sec:mpc-lqg} as the principal piece of remaining
analytical work. Stage~H of the project has since implemented it in
\texttt{smc2fc/control/lqg/} (linearisation by JAX-autodiff, backward
Runge--Kutta Riccati solver, \texttt{LQGController} emitting both
open-loop and feedback gains) and benchmarked it across the same
horizons as the SMC${}^2$ MPC. The empirical verdict reported in
Section~\ref{sec:results-lqg} is unambiguous, and we record it here
as the headline finding of this section.

\paragraph{The cubic Stuart--Landau non-linearity is essential.} No
choice of cost weights $(Q,R,Q_T,x_{\mathrm{ref}})$ within the LQG
family lets a Riccati-based controller match the SMC${}^2$ MPC on
the FSA-v2 plant. At the only horizon where both controllers
currently have committed checkpoints ($T = 14$ d), the SMC${}^2$ MPC
beats the tuned LQG by $15$ percentage points of the constant
baseline (ratio $1.157$ vs $1.008$); see
Figure~\ref{fig:t14}. The reason is structural and was diagnosed
in Section~\ref{sec:mpc-lqg}: the linearisation
\eqref{eq:lqg-numbers} has $B_{\mathrm{lin}}^{(A)} = 0$ (no direct
linear path from $\Phi$ to $A$), one positive eigenvalue ($+0.003$,
the cubic damping at $A_{\mathrm{typ}}$ is too weak to stabilise the
linearisation), and an inverted sign on the F$\to$A coupling
($A_{\mathrm{lin}}^{(A,F)} = -0.026$, opposite to the true plant).
At long horizons the slowly unstable eigenvalue dominates the
cost-to-go and the LQG schedule pushes $\Phi$ to bang-bang, breaking
the soft fatigue barrier (Table~\ref{tab:lqg-horizons}).

\paragraph{LQG is still useful as a baseline and as a warm-start
candidate.} Up to $T = 56$ d the tuned LQG outperforms the
constant-$\Phi=1$ baseline by $0$--$18\%$ with no fatigue-barrier
violation. It costs roughly five orders of magnitude less compute
than the SMC${}^2$ MPC (Section~\ref{sec:results-comparison}), so
it is a defensible cheap baseline against which any future SMC${}^2$
result should be measured. A natural follow-up -- still outstanding --
is to use the tuned LQG schedule as the prior mean for the
SMC${}^2$ controller's Gaussian prior over $\thetactrl$, with a
tightened standard deviation around it. Whether this meaningfully
reduces SMC${}^2$ compute at intermediate horizons
$T \in [28,56]$ d is an experimental question for a future
iteration.

\subsection{Other open questions}\label{sec:discussion-other}

A few smaller points deserve mention.

\begin{itemize}
  \item \emph{Horizon length and the SMC${}^2$ MPC sweep.} Stage~I has
        implemented the $1$-hour outer step (Section~\ref{sec:results-stage-i})
        recommended in Section~\ref{sec:num-stab-recommend} and run the
        full LQG sweep across
        $T \in \{14, 28, 42, 56, 84\}$ d. The corresponding SMC${}^2$
        MPC sweep is in progress on the math laptop. Once both controllers
        have committed checkpoints at every horizon, the
        SMC${}^2$-vs-LQG gap of Figure~\ref{fig:t14} will be measurable
        across the regime where LQG remains competitive ($T \le 56$ d)
        and the regime where it fails ($T = 84$ d), giving a clean
        quantification of the value added by the cubic Stuart--Landau
        non-linearity as a function of horizon.
  \item \emph{Posterior shape vs.\ posterior mean.} The controller of
        Section~\ref{sec:mpc-pipeline} uses only the posterior mean
        $\bar\theta_{\mathrm{ctrl}}$. The flat-cost-surface argument of
        Section~\ref{sec:mpc-robustness} explains why this is sometimes
        sufficient, but in regimes where the cost surface is curved the
        full posterior carries useful information. A risk-sensitive
        variant that minimises a quantile of the cost rather than the
        mean would be the natural extension.
  \item \emph{Parameter drift across windows.} The Bures--Wasserstein
        bridge of Section~\ref{sec:smc2} produces a sensible
        Gaussian-fit warm-start, but does not protect against
        \emph{slow} parameter drift in the truth (e.g.\ the athlete's
        time constants change with training age). A more principled
        approach would augment the static parameters with a slow
        random-walk component and let the filter track the drift
        explicitly.
\end{itemize}

\medskip
\noindent
With the FIM, Lyapunov, and LQG / Riccati analyses now all in place
-- and the empirical finding that no LQG variant matches the
SMC${}^2$ MPC on this plant -- the project has answered the three
analytical questions earlier drafts of these notes flagged as
outstanding. The model has empirical validation through the Stage~D,
E5 and H results, analytical guarantees on identifiability and
stability, and a quantitative case for the value added by the cubic
Stuart--Landau non-linearity. The remaining work is empirical: a
multi-horizon SMC${}^2$ MPC sweep against the LQG baseline, and the
LQG-warm-start variant of Section~\ref{sec:discussion-lqg}.
